\documentclass{article}
\usepackage[utf8]{inputenc}
\usepackage[russian]{babel}
\usepackage{amsmath, amssymb}
\usepackage{geometry}
\geometry{a4paper, margin=2cm}

\begin{document}


\section*{Задача I.6.3}

\textbf{Условие:} Пусть $y^*$ --- приближение к корню уравнения $f(y) = 0$. Вывести приближенное равенство $y - y^* \approx -\frac{f(y^*)}{f'(y^*)}$.

\vspace{0.5cm}
\textbf{Решение:}

Пусть $y$ --- точный корень уравнения $f(y) = 0$, а $y^*$ --- его некоторое приближенное значение. Предположим, что функция $f(y)$ является достаточно гладкой (непрерывно дифференцируемой) в окрестности точки $y^*$.

Разложим функцию $f(y)$ в ряд Тейлора в окрестности точки $y^*$:
\begin{equation}
f(y) = f(y^*) + f'(y^*)(y - y^*) + \frac{f''(\xi)}{2}(y - y^*)^2,
\end{equation}
где $\xi$ --- некоторая точка, лежащая между $y$ и $y^*$.

Так как $y$ является точным корнем уравнения, выполняется равенство $f(y) = 0$. Подставим его в разложение:
\begin{equation}
0 = f(y^*) + f'(y^*)(y - y^*) + \frac{f''(\xi)}{2}(y - y^*)^2.
\end{equation}

Поскольку $y^*$ выбрано как приближение к корню, разность $(y - y^*)$ является малой величиной. Следовательно, квадратом этой разности $(y - y^*)^2$ (членом второго порядка малости) можно пренебречь по сравнению с линейным членом. Отбрасывая остаточный член, получаем приближенное линейное соотношение:
\begin{equation}
0 \approx f(y^*) + f'(y^*)(y - y^*).
\end{equation}

Выразим из этого равенства искомую разность $y - y^*$:
\begin{equation}
f'(y^*)(y - y^*) \approx -f(y^*),
\end{equation}
\begin{equation}
y - y^* \approx -\frac{f(y^*)}{f'(y^*)}.
\end{equation}

\textbf{Замечание:} Полученное приближенное равенство имеет фундаментальное значение в вычислительной математике. Оно показывает, что поправка $\Delta y = y - y^*$, которую необходимо добавить к приближению $y^*$, чтобы получить точный корень, пропорциональна значению функции в этой точке и обратно пропорциональна её производной. Именно эта формула лежит в основе итерационного метода Ньютона (метода касательных) для численного решения нелинейных уравнений.





\section*{Задача I.8.11}

\textbf{Условие:} Найти линейную оценку погрешности для функции $u = \ln t$, если заданы точка приближения $t^* = 1$, значение функции $u^*$ в этой точке и погрешность $\Delta t^* = 0{,}1$.

\vspace{0.5cm}
\textbf{Решение:}

\subsection*{1. Исходные данные}

Задана функция:
\begin{equation}
u(t) = \ln t.
\end{equation}

Приближённое значение аргумента:
\begin{equation}
t^* = 1.
\end{equation}

Абсолютная погрешность аргумента:
\begin{equation}
\Delta t^* = 0{,}1.
\end{equation}

Значение функции в точке приближения:
\begin{equation}
u^* = u(t^*) = \ln 1 = 0.
\end{equation}

\subsection*{2. Формула линейной оценки погрешности}

Если функция $u(t)$ дифференцируема в окрестности точки $t^*$, то абсолютная погрешность функции $\Delta u^*$ оценивается по формуле:
\begin{equation}
\Delta u^* \approx \left| \frac{du}{dt}\bigg|_{t = t^*} \right| \cdot \Delta t^*.
\label{eq:lin_err}
\end{equation}

Эта формула следует из разложения функции в ряд Тейлора в окрестности точки $t^*$ с отбрасыванием членов второго и высших порядков малости (линеаризация).

\subsection*{3. Вычисление производной}

Найдём производную функции $u(t) = \ln t$:
\begin{equation}
\frac{du}{dt} = \frac{1}{t}.
\end{equation}

Вычислим значение производной в точке $t^* = 1$:
\begin{equation}
\frac{du}{dt}\bigg|_{t = 1} = \frac{1}{1} = 1.
\end{equation}

\subsection*{4. Линейная оценка абсолютной погрешности}

Подставим найденные значения в формулу~\eqref{eq:lin_err}:
\begin{equation}
\Delta u^* \approx |1| \cdot 0{,}1 = 0{,}1.
\end{equation}

Таким образом, линейная оценка абсолютной погрешности функции $u = \ln t$ при $t^* = 1$ и $\Delta t^* = 0{,}1$ равна:
\begin{equation}
\boxed{\Delta u^* \approx 0{,}1.}
\end{equation}

\subsection*{5. Относительная погрешность}

Относительная погрешность функции определяется как:
\begin{equation}
\delta u^* = \frac{\Delta u^*}{|u^*|}.
\end{equation}

Однако в нашем случае $u^* = \ln 1 = 0$, поэтому относительная погрешность \textbf{не определена} (формально стремится к бесконечности). Это характерная ситуация для функций, обращающихся в нуль в точке приближения: абсолютная погрешность может быть конечной, но относительная теряет смысл.

\subsection*{6. Проверка результата}

Для контроля вычислим точные значения функции на границах интервала неопределённости аргумента $t \in [t^* - \Delta t^*, t^* + \Delta t^*] = [0{,}9; 1{,}1]$:
\begin{align}
u(0{,}9) &= \ln 0{,}9 \approx -0{,}1054, \\
u(1{,}1) &= \ln 1{,}1 \approx 0{,}0953.
\end{align}

Максимальное отклонение от $u^* = 0$:
\begin{equation}
\max\bigl(|u(0{,}9) - 0|, |u(1{,}1) - 0|\bigr) = \max(0{,}1054; 0{,}0953) \approx 0{,}1054.
\end{equation}

Линейная оценка $\Delta u^* \approx 0{,}1$ близка к точному значению $0{,}1054$, что подтверждает корректность применения формулы~\eqref{eq:lin_err}. Небольшое расхождение объясняется отброшенными членами высшего порядка малости в разложении Тейлора.



\section*{Задача I.8.33}

\textbf{Условие:} Для вычисления первой производной функции $f(x)$ в точке $x + h$ используется формула $(f(x+2h) - f(x-2h))/(4h)$.

а) Каков порядок аппроксимации этой формулы?

б) Найти оптимальный шаг дифференцирования по этой формуле в произвольной точке для четырежды дифференцируемой функции.

в) Оценить его численное значение для функции $f(x) = \cos(x + \pi/4)$ в случае использования IEEE-арифметики одинарной и двойной точности.

\vspace{0.5cm}
\textbf{Решение:}

\subsection*{а) Порядок аппроксимации}

Нам нужно оценить погрешность формулы:
\begin{equation}
f'(x+h) \approx \frac{f(x+2h) - f(x-2h)}{4h}.
\label{eq:formula}
\end{equation}

Разложим $f(x+2h)$ и $f(x-2h)$ в ряд Тейлора в точке $x+h$. Заметим, что:
\begin{align}
x + 2h &= (x+h) + h, \\
x - 2h &= (x+h) - 3h.
\end{align}

Тогда разложения имеют вид:
\begin{align}
f(x+2h) &= f(x+h) + h f'(x+h) + \frac{h^2}{2} f''(x+h) + \frac{h^3}{6} f'''(x+h) + O(h^4), \\
f(x-2h) &= f(x+h) - 3h f'(x+h) + \frac{9h^2}{2} f''(x+h) - \frac{27h^3}{6} f'''(x+h) + O(h^4).
\end{align}

Вычислим разность:
\begin{align}
f(x+2h) - f(x-2h) &= 4h f'(x+h) - 4h^2 f''(x+h) + \frac{28h^3}{6} f'''(x+h) + O(h^4) \nonumber \\
&= 4h f'(x+h) - 4h^2 f''(x+h) + \frac{14h^3}{3} f'''(x+h) + O(h^4).
\end{align}

Подставим в формулу~\eqref{eq:formula}:
\begin{equation}
\frac{f(x+2h) - f(x-2h)}{4h} = f'(x+h) - h f''(x+h) + \frac{7h^2}{6} f'''(x+h) + O(h^3).
\end{equation}

Погрешность метода (усечение ряда Тейлора):
\begin{equation}
E_{\text{мет}} = \left| \frac{f(x+2h) - f(x-2h)}{4h} - f'(x+h) \right| = \left| -h f''(x+h) + O(h^2) \right| \approx h |f''(x+h)|.
\end{equation}

Поскольку главный член погрешности пропорционален $h^1$, \textbf{порядок аппроксимации формулы равен 1} (первый порядок).

\textbf{Замечание:} Это отличается от стандартной формулы центральной разности $(f(x+h) - f(x-h))/(2h)$, которая имеет второй порядок аппроксимации. В нашей формуле точки $x+2h$ и $x-2h$ симметричны относительно $x$, а не относительно $x+h$, что приводит к появлению члена с $f''(x+h)$ и снижению порядка точности.

\subsection*{б) Оптимальный шаг дифференцирования}

Полная погрешность численного дифференцирования складывается из двух составляющих:
\begin{enumerate}
\item \textbf{Погрешность метода} (усечение ряда Тейлора): $E_{\text{мет}} \approx h |f''(x+h)|$.
\item \textbf{Погрешность вычислений} (ошибки округления при вычислении $f$): при вычислении разности $f(x+2h) - f(x-2h)$ каждое значение функции вычисляется с относительной погрешностью $\varepsilon_{\text{маш}}$. Абсолютная погрешность разности порядка $2\varepsilon_{\text{маш}} |f|$, а после деления на $4h$ получаем:
\begin{equation}
E_{\text{выч}} \approx \frac{2\varepsilon_{\text{маш}} |f(x+h)|}{4h} = \frac{\varepsilon_{\text{маш}} |f(x+h)|}{2h}.
\end{equation}
\end{enumerate}

Полная погрешность:
\begin{equation}
E_{\text{полн}} \approx h |f''(x+h)| + \frac{\varepsilon_{\text{маш}} |f(x+h)|}{2h}.
\label{eq:total_error}
\end{equation}

Найдём оптимальный шаг $h_{\text{opt}}$, минимизирующий $E_{\text{полн}}$. Продифференцируем~\eqref{eq:total_error} по $h$ и приравняем к нулю:
\begin{equation}
\frac{dE_{\text{полн}}}{dh} = |f''(x+h)| - \frac{\varepsilon_{\text{маш}} |f(x+h)|}{2h^2} = 0.
\end{equation}

Отсюда:
\begin{equation}
h^2 = \frac{\varepsilon_{\text{маш}} |f(x+h)|}{2|f''(x+h)|},
\end{equation}
\begin{equation}
\boxed{h_{\text{opt}} = \sqrt{\frac{\varepsilon_{\text{маш}} |f(x+h)|}{2|f''(x+h)|}}.}
\label{eq:h_opt}
\end{equation}

\textbf{Замечание:} Сравните с формулой оптимального шага для центральной разности второго порядка: там $h_{\text{opt}} \sim \varepsilon_{\text{маш}}^{1/3}$, а у нас $h_{\text{opt}} \sim \varepsilon_{\text{маш}}^{1/2}$. Это следствие более низкого порядка аппроксимации.

\subsection*{в) Численная оценка для $f(x) = \cos(x + \pi/4)$}

Найдём производные:
\begin{align}
f(x) &= \cos(x + \pi/4), \\
f'(x) &= -\sin(x + \pi/4), \\
f''(x) &= -\cos(x + \pi/4).
\end{align}

В точке $x+h$:
\begin{align}
f(x+h) &= \cos(x+h+\pi/4), \\
f''(x+h) &= -\cos(x+h+\pi/4).
\end{align}

Отношение:
\begin{equation}
\frac{|f(x+h)|}{|f''(x+h)|} = \frac{|\cos(x+h+\pi/4)|}{|\cos(x+h+\pi/4)|} = 1 \quad (\text{при } \cos(x+h+\pi/4) \neq 0).
\end{equation}

Подставим в формулу~\eqref{eq:h_opt}:
\begin{equation}
h_{\text{opt}} = \sqrt{\frac{\varepsilon_{\text{маш}}}{2}}.
\end{equation}

\textbf{Для IEEE-арифметики одинарной точности} ($\varepsilon_{\text{маш}} \approx 6{,}0 \cdot 10^{-8}$):
\begin{equation}
h_{\text{opt}}^{(\text{float})} = \sqrt{\frac{6{,}0 \cdot 10^{-8}}{2}} = \sqrt{3{,}0 \cdot 10^{-8}} \approx 1{,}7 \cdot 10^{-4}.
\end{equation}

\textbf{Для IEEE-арифметики двойной точности} ($\varepsilon_{\text{маш}} \approx 1{,}1 \cdot 10^{-16}$):
\begin{equation}
h_{\text{opt}}^{(\text{double})} = \sqrt{\frac{1{,}1 \cdot 10^{-16}}{2}} = \sqrt{5{,}5 \cdot 10^{-17}} \approx 7{,}4 \cdot 10^{-9}.
\end{equation}

\subsection*{Итог}

\begin{itemize}
\item Порядок аппроксимации формулы $(f(x+2h) - f(x-2h))/(4h)$ для вычисления $f'(x+h)$ равен \textbf{1} (первый порядок).
\item Оптимальный шаг: $h_{\text{opt}} = \sqrt{\varepsilon_{\text{маш}} |f| / (2|f''|)}$.
\item Для $f(x) = \cos(x + \pi/4)$: $h_{\text{opt}} \approx 1{,}7 \cdot 10^{-4}$ (одинарная точность) и $h_{\text{opt}} \approx 7{,}4 \cdot 10^{-9}$ (двойная точность).
\end{itemize}






\section*{Задача II.7.39}

\textbf{Условие:} Доказать, что для систем линейных уравнений второго порядка ($n = 2$) методы Якоби и Гаусса--Зейделя сходятся и расходятся одновременно.

\vspace{0.5cm}
\textbf{Решение:}

\subsection*{1. Постановка задачи}

Рассмотрим систему линейных уравнений $\mathbf{A}\mathbf{x} = \mathbf{b}$, где $\mathbf{A}$ --- матрица размера $2 \times 2$:
\begin{equation}
\mathbf{A} =
\begin{pmatrix}
a_{11} & a_{12} \\
a_{21} & a_{22}
\end{pmatrix}, \quad
\mathbf{x} =
\begin{pmatrix}
x_1 \\ x_2
\end{pmatrix}, \quad
\mathbf{b} =
\begin{pmatrix}
b_1 \\ b_2
\end{pmatrix}.
\end{equation}

Предполагаем, что диагональные элементы ненулевые: $a_{11} \neq 0$, $a_{22} \neq 0$.

Разложим матрицу $\mathbf{A}$ в виде:
\begin{equation}
\mathbf{A} = \mathbf{D} - \mathbf{L} - \mathbf{U},
\end{equation}
где $\mathbf{D}$ --- диагональная часть, $-\mathbf{L}$ --- строго нижняя треугольная, $-\mathbf{U}$ --- строго верхняя треугольная:
\begin{equation}
\mathbf{D} =
\begin{pmatrix}
a_{11} & 0 \\
0 & a_{22}
\end{pmatrix}, \quad
\mathbf{L} =
\begin{pmatrix}
0 & 0 \\
-a_{21} & 0
\end{pmatrix}, \quad
\mathbf{U} =
\begin{pmatrix}
0 & -a_{12} \\
0 & 0
\end{pmatrix}.
\end{equation}

\subsection*{2. Метод Якоби}

Итерационная матрица метода Якоби имеет вид:
\begin{equation}
\mathbf{B}_J = \mathbf{D}^{-1}(\mathbf{L} + \mathbf{U}).
\end{equation}

Вычислим её явно:
\begin{equation}
\mathbf{B}_J =
\begin{pmatrix}
1/a_{11} & 0 \\
0 & 1/a_{22}
\end{pmatrix}
\begin{pmatrix}
0 & -a_{12} \\
-a_{21} & 0
\end{pmatrix}
=
\begin{pmatrix}
0 & -a_{12}/a_{11} \\
-a_{21}/a_{22} & 0
\end{pmatrix}.
\end{equation}

Найдём собственные значения $\mathbf{B}_J$ из характеристического уравнения $\det(\mathbf{B}_J - \lambda \mathbf{I}) = 0$:
\begin{equation}
\det
\begin{pmatrix}
-\lambda & -a_{12}/a_{11} \\
-a_{21}/a_{22} & -\lambda
\end{pmatrix}
= \lambda^2 - \frac{a_{12} a_{21}}{a_{11} a_{22}} = 0.
\end{equation}

Отсюда собственные значения:
\begin{equation}
\lambda_{1,2}^{(J)} = \pm \sqrt{\frac{a_{12} a_{21}}{a_{11} a_{22}}}.
\end{equation}

Спектральный радиус матрицы Якоби:
\begin{equation}
\rho(\mathbf{B}_J) = \sqrt{\left|\frac{a_{12} a_{21}}{a_{11} a_{22}}\right|}.
\label{eq:rho_j}
\end{equation}

Метод Якоби сходится тогда и только тогда, когда $\rho(\mathbf{B}_J) < 1$.

\subsection*{3. Метод Гаусса--Зейделя}

Итерационная матрица метода Гаусса--Зейделя:
\begin{equation}
\mathbf{B}_{GS} = (\mathbf{D} - \mathbf{L})^{-1} \mathbf{U}.
\end{equation}

Сначала найдём $(\mathbf{D} - \mathbf{L})^{-1}$:
\begin{equation}
\mathbf{D} - \mathbf{L} =
\begin{pmatrix}
a_{11} & 0 \\
a_{21} & a_{22}
\end{pmatrix}, \quad
(\mathbf{D} - \mathbf{L})^{-1} =
\begin{pmatrix}
1/a_{11} & 0 \\
-a_{21}/(a_{11} a_{22}) & 1/a_{22}
\end{pmatrix}.
\end{equation}

Теперь вычислим $\mathbf{B}_{GS}$:
\begin{equation}
\mathbf{B}_{GS} =
\begin{pmatrix}
1/a_{11} & 0 \\
-a_{21}/(a_{11} a_{22}) & 1/a_{22}
\end{pmatrix}
\begin{pmatrix}
0 & -a_{12} \\
0 & 0
\end{pmatrix}
=
\begin{pmatrix}
0 & -a_{12}/a_{11} \\
0 & a_{12} a_{21}/(a_{11} a_{22})
\end{pmatrix}.
\end{equation}

Матрица $\mathbf{B}_{GS}$ является верхнетреугольной, поэтому её собственные значения равны диагональным элементам:
\begin{equation}
\lambda_1^{(GS)} = 0, \quad \lambda_2^{(GS)} = \frac{a_{12} a_{21}}{a_{11} a_{22}}.
\end{equation}

Спектральный радиус матрицы Гаусса--Зейделя:
\begin{equation}
\rho(\mathbf{B}_{GS}) = \left|\frac{a_{12} a_{21}}{a_{11} a_{22}}\right|.
\label{eq:rho_gs}
\end{equation}

\subsection*{4. Сравнение спектральных радиусов}

Сравнивая формулы~\eqref{eq:rho_j} и~\eqref{eq:rho_gs}, получаем фундаментальное соотношение для систем второго порядка:
\begin{equation}
\boxed{\rho(\mathbf{B}_{GS}) = \left[\rho(\mathbf{B}_J)\right]^2.}
\label{eq:key_relation}
\end{equation}

\subsection*{5. Доказательство одновременности сходимости и расходимости}

Из соотношения~\eqref{eq:key_relation} непосредственно следуют оба утверждения:

\textbf{Сходимость:}
\begin{equation}
\rho(\mathbf{B}_J) < 1 \iff \left[\rho(\mathbf{B}_J)\right]^2 < 1 \iff \rho(\mathbf{B}_{GS}) < 1.
\end{equation}
То есть метод Якоби сходится тогда и только тогда, когда сходится метод Гаусса--Зейделя.

\textbf{Расходимость:}
\begin{equation}
\rho(\mathbf{B}_J) > 1 \iff \left[\rho(\mathbf{B}_J)\right]^2 > 1 \iff \rho(\mathbf{B}_{GS}) > 1.
\end{equation}
То есть метод Якоби расходится тогда и только тогда, когда расходится метод Гаусса--Зейделя.

\textbf{Граничный случай} $\rho(\mathbf{B}_J) = 1$ также совпадает: $\rho(\mathbf{B}_{GS}) = 1$.

\subsection*{6. Замечание о скорости сходимости}

Из соотношения $\rho(\mathbf{B}_{GS}) = [\rho(\mathbf{B}_J)]^2$ следует важный практический вывод: если оба метода сходятся, то метод Гаусса--Зейделя сходится \textbf{примерно в два раза быстрее} метода Якоби (поскольку число итераций, необходимое для достижения заданной точности, обратно пропорционально $-\ln \rho$).

\textbf{Важно:} данное свойство (одновременность сходимости) является специфическим для систем размера $n = 2$. Для $n \geqslant 3$ существуют примеры, когда один из методов сходится, а другой --- расходится.



\section*{Задача II.9.2 (пункт б)}

\textbf{Условие:} Дана система линейных алгебраических уравнений $\mathbf{A}\mathbf{x} = \mathbf{b}$. Оценить максимально точную относительную погрешность $\|\Delta\mathbf{x}\|/\|\mathbf{x}\|$ в заданной норме. Найти вектор ошибки $\Delta\mathbf{b}$, на котором эта оценка достигается. При каком $\Delta\mathbf{b}$ относительная ошибка $\|\Delta\mathbf{x}\|/\|\mathbf{x}\|$ будет минимальной? Найти её.

$$\mathbf{A} = \begin{pmatrix} 101 & -110 \\ -110 & 122 \end{pmatrix}, \quad \mathbf{b} = \begin{pmatrix} 92 \\ -98 \end{pmatrix}, \quad \frac{\|\Delta\mathbf{b}\|}{\|\mathbf{b}\|} = 0{,}01, \quad \|\mathbf{x}\| = \sum_i |x_i|.$$

\textit{Примечание:} в условии норма обозначена как $\|\cdot\|_2$, но по формуле $\sum_i |x_i|$ это стандартная норма $\|\cdot\|_1$ (сумма модулей компонент). Будем следовать обозначению условия, но помнить, что работаем с суммой модулей.

\vspace{0.5cm}
\textbf{Решение:}

\subsection*{1. Решение системы $\mathbf{A}\mathbf{x} = \mathbf{b}$}

Найдём определитель матрицы $\mathbf{A}$:
\begin{equation}
\det(\mathbf{A}) = 101 \cdot 122 - (-110) \cdot (-110) = 12322 - 12100 = 222.
\end{equation}

Обратная матрица:
\begin{equation}
\mathbf{A}^{-1} = \frac{1}{222} \begin{pmatrix} 122 & 110 \\ 110 & 101 \end{pmatrix}.
\end{equation}

Решение системы $\mathbf{x} = \mathbf{A}^{-1}\mathbf{b}$:
\begin{equation}
\mathbf{x} = \frac{1}{222} \begin{pmatrix} 122 & 110 \\ 110 & 101 \end{pmatrix} \begin{pmatrix} 92 \\ -98 \end{pmatrix} = \frac{1}{222} \begin{pmatrix} 11224 - 10780 \\ 10120 - 9898 \end{pmatrix} = \frac{1}{222} \begin{pmatrix} 444 \\ 222 \end{pmatrix} = \begin{pmatrix} 2 \\ 1 \end{pmatrix}.
\end{equation}

Итак, $\mathbf{x} = (2, 1)^T$.

\subsection*{2. Вычисление норм}

Норма вектора $\mathbf{x}$ (сумма модулей компонент):
\begin{equation}
\|\mathbf{x}\| = |2| + |1| = 3.
\end{equation}

Норма вектора $\mathbf{b}$:
\begin{equation}
\|\mathbf{b}\| = |92| + |-98| = 190.
\end{equation}

Абсолютная погрешность правой части:
\begin{equation}
\|\Delta\mathbf{b}\| = 0{,}01 \cdot \|\mathbf{b}\| = 0{,}01 \cdot 190 = 1{,}9.
\end{equation}

Норма матрицы $\mathbf{A}$ (максимум сумм модулей элементов по столбцам):
\begin{equation}
\|\mathbf{A}\| = \max\bigl(|101| + |-110|,\; |-110| + |122|\bigr) = \max(211,\, 232) = 232.
\end{equation}

Норма обратной матрицы:
\begin{equation}
\|\mathbf{A}^{-1}\| = \frac{1}{222} \max\bigl(|122| + |110|,\; |110| + |101|\bigr) = \frac{1}{222} \max(232,\, 211) = \frac{232}{222} = \frac{116}{111}.
\end{equation}

\subsection*{3. Число обусловленности и теоретическая оценка погрешности}

Число обусловленности матрицы $\mathbf{A}$:
\begin{equation}
\operatorname{cond}(\mathbf{A}) = \|\mathbf{A}\| \cdot \|\mathbf{A}^{-1}\| = 232 \cdot \frac{232}{222} = \frac{53824}{222} = \frac{26912}{111} \approx 242{,}45.
\end{equation}

\textbf{Важное замечание о смысле этой оценки.} Стандартная теорема даёт \textit{верхнюю границу} относительной погрешности:
\begin{equation}
\frac{\|\Delta\mathbf{x}\|}{\|\mathbf{x}\|} \leqslant \operatorname{cond}(\mathbf{A}) \cdot \frac{\|\Delta\mathbf{b}\|}{\|\mathbf{b}\|} = \frac{26912}{111} \cdot 0{,}01 \approx 2{,}4245.
\end{equation}
Эта оценка \textbf{достижима в худшем случае} --- когда одновременно и $\mathbf{b}$, и $\Delta\mathbf{b}$ направлены ``наихудшим'' образом. Однако в нашей задаче вектор $\mathbf{b}$ \textit{фиксирован}, поэтому реальная максимальная погрешность будет меньше. Найдём её точно.

\subsection*{4. Точный максимум погрешности при фиксированном $\mathbf{b}$}

Из уравнения $\mathbf{A}(\mathbf{x} + \Delta\mathbf{x}) = \mathbf{b} + \Delta\mathbf{b}$ следует $\Delta\mathbf{x} = \mathbf{A}^{-1}\Delta\mathbf{b}$. Поскольку $\mathbf{x}$ фиксирован, задача сводится к максимизации $\|\mathbf{A}^{-1}\Delta\mathbf{b}\|$ при условии $\|\Delta\mathbf{b}\| = 1{,}9$.

\textbf{Как найти максимум $\|\mathbf{A}^{-1}\Delta\mathbf{b}\|$?} В норме $\|\cdot\| = \sum |x_i|$ максимум $\|\mathbf{M}\mathbf{y}\|/\|\mathbf{y}\|$ равен $\|\mathbf{M}\|$ и достигается, когда $\mathbf{y}$ --- это базисный вектор $\mathbf{e}_j$, соответствующий столбцу матрицы $\mathbf{M}$ с наибольшей суммой модулей элементов.

Применяем к $\mathbf{M} = \mathbf{A}^{-1}$:
\begin{itemize}
    \item Первый столбец $\mathbf{A}^{-1}$: сумма модулей $= (122 + 110)/222 = 232/222$.
    \item Второй столбец $\mathbf{A}^{-1}$: сумма модулей $= (110 + 101)/222 = 211/222$.
\end{itemize}
Первый столбец даёт максимум, поэтому оптимальное направление: $\Delta\mathbf{b} = \pm 1{,}9 \cdot \mathbf{e}_1 = \pm(1{,}9;\, 0)^T$.

\textbf{Проверка.} При $\Delta\mathbf{b} = (1{,}9;\, 0)^T$:
\begin{equation}
\Delta\mathbf{x} = \mathbf{A}^{-1}\Delta\mathbf{b} = \frac{1{,}9}{222} \begin{pmatrix} 122 \\ 110 \end{pmatrix} \approx \begin{pmatrix} 1{,}044 \\ 0{,}941 \end{pmatrix}.
\end{equation}
\begin{equation}
\|\Delta\mathbf{x}\| = \frac{1{,}9}{222}(122 + 110) = \frac{1{,}9 \cdot 232}{222} \approx 1{,}986.
\end{equation}
\begin{equation}
\frac{\|\Delta\mathbf{x}\|}{\|\mathbf{x}\|} = \frac{1{,}986}{3} \approx 0{,}662.
\end{equation}

\textbf{Итог по максимуму:}
\begin{equation}
\boxed{\max_{\|\Delta\mathbf{b}\|=1{,}9} \frac{\|\Delta\mathbf{x}\|}{\|\mathbf{x}\|} = \frac{\|\mathbf{A}^{-1}\| \cdot \|\Delta\mathbf{b}\|}{\|\mathbf{x}\|} = \frac{232 \cdot 1{,}9}{222 \cdot 3} = \frac{440{,}8}{666} \approx 0{,}662.}
\end{equation}
Оценка достигается на $\Delta\mathbf{b} = \pm(1{,}9;\, 0)^T$.

\textbf{Почему оценка через $\operatorname{cond}(\mathbf{A}) \approx 2{,}42$ не достигается?} Потому что она предполагает ``худший'' $\mathbf{b}$, а наш $\mathbf{b} = (92;\, -98)^T$ не является таким. Отношение $\|\mathbf{b}\|/\|\mathbf{x}\| = 190/3 \approx 63{,}3$, тогда как для ``худшего'' $\mathbf{b}$ это отношение было бы равно $\|\mathbf{A}\| = 232$.

\subsection*{5. Минимальная погрешность}

\textbf{Тривиальный случай:} при $\Delta\mathbf{b} = \mathbf{0}$ погрешность равна нулю:
\begin{equation}
\frac{\|\Delta\mathbf{x}\|}{\|\mathbf{x}\|} = 0.
\end{equation}

\textbf{Нетривиальный случай:} найдём минимум $\|\Delta\mathbf{x}\|/\|\mathbf{x}\|$ при фиксированном $\|\Delta\mathbf{b}\| = 1{,}9$, т.\,е. минимум $\|\mathbf{A}^{-1}\Delta\mathbf{b}\|/\|\Delta\mathbf{b}\|$.

\textbf{Ключевое свойство подчинённых норм:} для любой обратимой матрицы $\mathbf{M}$
\begin{equation}
\min_{\mathbf{y} \neq 0} \frac{\|\mathbf{M}\mathbf{y}\|}{\|\mathbf{y}\|} = \frac{1}{\|\mathbf{M}^{-1}\|}.
\end{equation}
\textit{Доказательство:} из $\mathbf{y} = \mathbf{M}^{-1}(\mathbf{M}\mathbf{y})$ и неравенства $\|\mathbf{y}\| \leqslant \|\mathbf{M}^{-1}\| \cdot \|\mathbf{M}\mathbf{y}\|$ следует $\|\mathbf{M}\mathbf{y}\|/\|\mathbf{y}\| \geqslant 1/\|\mathbf{M}^{-1}\|$. Равенство достигается на ``худшем'' векторе для $\mathbf{M}^{-1}$.

Применяем к $\mathbf{M} = \mathbf{A}^{-1}$:
\begin{equation}
\min \frac{\|\mathbf{A}^{-1}\Delta\mathbf{b}\|}{\|\Delta\mathbf{b}\|} = \frac{1}{\|(\mathbf{A}^{-1})^{-1}\|} = \frac{1}{\|\mathbf{A}\|} = \frac{1}{232}.
\end{equation}

\textbf{На каком $\Delta\mathbf{b}$ достигается минимум?} На векторе, где $\mathbf{A}$ достигает своей нормы, т.\,е. на столбце $\mathbf{A}$ с максимальной суммой модулей элементов. Это второй столбец: $(-110;\, 122)^T$ (сумма модулей $= 232 = \|\mathbf{A}\|$).

Нормируем: $\|\Delta\mathbf{b}\| = 1{,}9$, поэтому
\begin{equation}
\Delta\mathbf{b}_{\min} = \frac{1{,}9}{232} \begin{pmatrix} -110 \\ 122 \end{pmatrix} \approx \begin{pmatrix} -0{,}901 \\ 1{,}000 \end{pmatrix}.
\end{equation}

\textbf{Проверка.} Вычислим $\mathbf{A}^{-1}\Delta\mathbf{b}_{\min}$:
\begin{equation}
\mathbf{A}^{-1} \begin{pmatrix} -110 \\ 122 \end{pmatrix} = \frac{1}{222} \begin{pmatrix} 122 \cdot (-110) + 110 \cdot 122 \\ 110 \cdot (-110) + 101 \cdot 122 \end{pmatrix} = \frac{1}{222} \begin{pmatrix} 0 \\ 222 \end{pmatrix} = \begin{pmatrix} 0 \\ 1 \end{pmatrix}.
\end{equation}
Значит,
\begin{equation}
\Delta\mathbf{x} = \frac{1{,}9}{232} \begin{pmatrix} 0 \\ 1 \end{pmatrix}, \quad \|\Delta\mathbf{x}\| = \frac{1{,}9}{232} \approx 0{,}00819.
\end{equation}
\begin{equation}
\frac{\|\Delta\mathbf{x}\|}{\|\mathbf{x}\|} = \frac{1{,}9}{232 \cdot 3} = \frac{1{,}9}{696} \approx 0{,}00273.
\end{equation}

\textbf{Итог по минимуму:}
\begin{equation}
\boxed{\min_{\|\Delta\mathbf{b}\|=1{,}9} \frac{\|\Delta\mathbf{x}\|}{\|\mathbf{x}\|} = \frac{\|\Delta\mathbf{b}\|}{\|\mathbf{A}\| \cdot \|\mathbf{x}\|} = \frac{1{,}9}{232 \cdot 3} \approx 0{,}00273 \quad (\approx 0{,}27\%).}
\end{equation}
Минимум достигается на $\Delta\mathbf{b} \propto (-110;\, 122)^T$.

\subsection*{Итог}

\begin{itemize}
    \item Решение системы: $\mathbf{x} = (2,\, 1)^T$.
    \item Теоретическая оценка через $\operatorname{cond}(\mathbf{A})$: $\|\Delta\mathbf{x}\|/\|\mathbf{x}\| \leqslant 2{,}42$ ($\approx 242\%$) --- достижима только при ``худшем'' $\mathbf{b}$.
    \item Реальный максимум при данном $\mathbf{b}$: $\approx 0{,}662$ ($\approx 66\%$), достигается на $\Delta\mathbf{b} = \pm(1{,}9;\, 0)^T$.
    \item Минимальная ненулевая погрешность: $\approx 0{,}00273$ ($\approx 0{,}27\%$), достигается на $\Delta\mathbf{b} \propto (-110;\, 122)^T$.
\end{itemize}
\textbf{Условие:} Дана система линейных алгебраических уравнений $\mathbf{A}\mathbf{x} = \mathbf{b}$. Оценить максимально точную относительную погрешность $\|\Delta\mathbf{x}\|/\|\mathbf{x}\|$ в заданной норме. Найти вектор ошибки $\Delta\mathbf{b}$, на котором эта оценка достигается. При каком $\Delta\mathbf{b}$ относительная ошибка $\|\Delta\mathbf{x}\|/\|\mathbf{x}\|$ будет минимальной? Найти её.

$$\mathbf{A} = \begin{pmatrix} 101 & -110 \\ -110 & 122 \end{pmatrix}, \quad \mathbf{b} = \begin{pmatrix} 92 \\ -98 \end{pmatrix}, \quad \frac{\|\Delta\mathbf{b}\|_2}{\|\mathbf{b}\|_2} = 0{,}01, \quad \|\mathbf{x}\|_2 = \sum_i |x_i|.$$

\vspace{0.5cm}
\textbf{Решение:}

\subsection*{1. Решение системы $\mathbf{A}\mathbf{x} = \mathbf{b}$}

Найдём определитель матрицы $\mathbf{A}$:
\begin{equation}
\det(\mathbf{A}) = 101 \cdot 122 - (-110) \cdot (-110) = 12322 - 12100 = 222.
\end{equation}

Обратная матрица:
\begin{equation}
\mathbf{A}^{-1} = \frac{1}{222} \begin{pmatrix} 122 & 110 \\ 110 & 101 \end{pmatrix}.
\end{equation}

Решение системы $\mathbf{x} = \mathbf{A}^{-1}\mathbf{b}$:
\begin{equation}
\mathbf{x} = \frac{1}{222} \begin{pmatrix} 122 & 110 \\ 110 & 101 \end{pmatrix} \begin{pmatrix} 92 \\ -98 \end{pmatrix} = \frac{1}{222} \begin{pmatrix} 122 \cdot 92 + 110 \cdot (-98) \\ 110 \cdot 92 + 101 \cdot (-98) \end{pmatrix} = \frac{1}{222} \begin{pmatrix} 11224 - 10780 \\ 10120 - 9898 \end{pmatrix} = \frac{1}{222} \begin{pmatrix} 444 \\ 222 \end{pmatrix} = \begin{pmatrix} 2 \\ 1 \end{pmatrix}.
\end{equation}

Итак, $\mathbf{x} = (2, 1)^T$.

\subsection*{2. Вычисление норм}

Согласно условию, используется норма $\|\mathbf{x}\|_2 = \sum_i |x_i|$ (сумма модулей компонент):
\begin{equation}
\|\mathbf{x}\|_2 = |2| + |1| = 3.
\end{equation}

Норма вектора $\mathbf{b}$:
\begin{equation}
\|\mathbf{b}\|_2 = |92| + |-98| = 190.
\end{equation}

Норма матрицы $\mathbf{A}$ (в той же норме, что и векторы --- максимум сумм модулей элементов по столбцам):
\begin{equation}
\|\mathbf{A}\|_2 = \max_j \sum_i |a_{ij}| = \max(|101| + |-110|, |-110| + |122|) = \max(211, 232) = 232.
\end{equation}

Норма обратной матрицы:
\begin{equation}
\|\mathbf{A}^{-1}\|_2 = \frac{1}{222} \max(|122| + |110|, |110| + |101|) = \frac{1}{222} \max(232, 211) = \frac{232}{222} = \frac{116}{111}.
\end{equation}

\subsection*{3. Число обусловленности и оценка погрешности}

Число обусловленности матрицы $\mathbf{A}$:
\begin{equation}
\text{cond}(\mathbf{A}) = \|\mathbf{A}\|_2 \cdot \|\mathbf{A}^{-1}\|_2 = 232 \cdot \frac{232}{222} = \frac{53824}{222} = \frac{26912}{111} \approx 242{,}45.
\end{equation}

Максимально точная оценка относительной погрешности:
\begin{equation}
\frac{\|\Delta\mathbf{x}\|_2}{\|\mathbf{x}\|_2} \leqslant \text{cond}(\mathbf{A}) \cdot \frac{\|\Delta\mathbf{b}\|_2}{\|\mathbf{b}\|_2} = \frac{26912}{111} \cdot 0{,}01 = \frac{269{,}12}{111} \approx 2{,}4245.
\end{equation}

Таким образом, относительная погрешность решения может достигать $\approx 242\%$, что свидетельствует о плохой обусловленности системы.

\subsection*{4. Вектор $\Delta\mathbf{b}$, на котором оценка достигается}

Оценка достигается, когда $\Delta\mathbf{b}$ направлен так, чтобы максимизировать $\|\mathbf{A}^{-1}\Delta\mathbf{b}\|_2$ при заданном $\|\Delta\mathbf{b}\|_2$.

Для нормы $\|\cdot\|_2 = \sum |x_i|$ максимум $\|\mathbf{A}^{-1}\Delta\mathbf{b}\|_2$ достигается, когда $\Delta\mathbf{b}$ пропорционален столбцу матрицы $\mathbf{A}^{-1}$ с максимальной суммой модулей элементов.

Из вычислений выше, второй столбец $\mathbf{A}^{-1}$ имеет большую сумму: $(110 + 101)/222 = 211/222$, а первый: $(122 + 110)/222 = 232/222$.

Следовательно, максимум достигается на первом столбце:
\begin{equation}
\Delta\mathbf{b} = \pm \|\Delta\mathbf{b}\|_2 \cdot \mathbf{e}_1 = \pm 0{,}01 \cdot 190 \cdot \begin{pmatrix} 1 \\ 0 \end{pmatrix} = \pm \begin{pmatrix} 1{,}9 \\ 0 \end{pmatrix}.
\end{equation}

Проверка: при $\Delta\mathbf{b} = (1{,}9, 0)^T$:
\begin{equation}
\Delta\mathbf{x} = \mathbf{A}^{-1}\Delta\mathbf{b} = \frac{1}{222} \begin{pmatrix} 122 & 110 \\ 110 & 101 \end{pmatrix} \begin{pmatrix} 1{,}9 \\ 0 \end{pmatrix} = \frac{1}{222} \begin{pmatrix} 231{,}8 \\ 209 \end{pmatrix} \approx \begin{pmatrix} 1{,}044 \\ 0{,}941 \end{pmatrix}.
\end{equation}

\begin{equation}
\|\Delta\mathbf{x}\|_2 = 1{,}044 + 0{,}941 = 1{,}985.
\end{equation}

\begin{equation}
\frac{\|\Delta\mathbf{x}\|_2}{\|\mathbf{x}\|_2} = \frac{1{,}985}{3} \approx 0{,}662.
\end{equation}

Это меньше теоретического максимума $2{,}4245$, потому что точное достижение оценки требует также оптимального направления $\mathbf{x}$, которое в нашей задаче фиксировано.

\subsection*{5. Минимальная относительная ошибка}

Минимальная относительная ошибка $\|\Delta\mathbf{x}\|_2/\|\mathbf{x}\|_2$ достигается при $\Delta\mathbf{b} = \mathbf{0}$ и равна:
\begin{equation}
\boxed{\frac{\|\Delta\mathbf{x}\|_2}{\|\mathbf{x}\|_2} = 0.}
\end{equation}

Если же требуется найти минимальную ненулевую ошибку при фиксированном $\|\Delta\mathbf{b}\|_2/\|\mathbf{b}\|_2 = 0{,}01$, то она достигается, когда $\Delta\mathbf{b}$ направлен в сторону, где $\mathbf{A}^{-1}$ имеет минимальное ``усиление''.

Для нормы $\|\cdot\|_2 = \sum |x_i|$ минимум $\|\mathbf{A}^{-1}\Delta\mathbf{b}\|_2/\|\Delta\mathbf{b}\|_2$ равен $1/\|\mathbf{A}\|_2$ (по свойству согласованных норм).

Следовательно:
\begin{equation}
\min \frac{\|\Delta\mathbf{x}\|_2}{\|\mathbf{x}\|_2} = \frac{1}{\|\mathbf{A}\|_2} \cdot \frac{\|\Delta\mathbf{b}\|_2}{\|\mathbf{b}\|_2} \cdot \frac{\|\mathbf{b}\|_2}{\|\mathbf{x}\|_2} = \frac{1}{232} \cdot 0{,}01 \cdot \frac{190}{3} \approx 0{,}00273.
\end{equation}

Это значение достигается, когда $\Delta\mathbf{b}$ пропорционален вектору, на котором $\mathbf{A}$ достигает своей нормы.

\subsection*{Итог}

\begin{itemize}
\item Решение системы: $\mathbf{x} = (2, 1)^T$.
\item Максимальная относительная погрешность: $\|\Delta\mathbf{x}\|_2/\|\mathbf{x}\|_2 \leqslant 2{,}4245$ ($\approx 242\%$).
\item Оценка достигается на $\Delta\mathbf{b} = \pm(1{,}9, 0)^T$ (с точностью до направления).
\item Минимальная относительная ошибка: $0$ (при $\Delta\mathbf{b} = \mathbf{0}$).
\item Минимальная ненулевая ошибка при заданном $\|\Delta\mathbf{b}\|_2$: $\approx 0{,}00273$ ($\approx 0{,}27\%$).
\end{itemize}



\section*{Задача II.9.14}

\textbf{Условие:} Система уравнений $\mathbf{A}\mathbf{x} = \mathbf{f}$, где
$$\mathbf{A} = \begin{pmatrix} 3 & -1 & 1 \\ 2 & 6 & -2 \\ -1 & 1 & 9 \end{pmatrix}, \quad \mathbf{f} = (4, -4, 2)^T,$$
решается с помощью метода простых итераций, начальное приближение нулевое. При каких значениях итерационного параметра метод будет сходиться:

а) при вычислениях с бесконечным числом бит в мантиссе?

б) при длине мантиссы 52 бита?

в) при каком значении итерационного параметра сходимость будет самая быстрая?

г) Найти значение оптимального параметра для первого шага решения системы методом наискорейшего спуска.

\vspace{0.5cm}
\textbf{Решение:}

\subsection*{Предварительный анализ: собственные значения матрицы $\mathbf{A}$}

Метод простых итераций с параметром $\tau$ имеет вид:
\begin{equation}
\mathbf{x}^{(k+1)} = \mathbf{x}^{(k)} - \tau(\mathbf{A}\mathbf{x}^{(k)} - \mathbf{f}) = (\mathbf{I} - \tau\mathbf{A})\mathbf{x}^{(k)} + \tau\mathbf{f}.
\end{equation}

Итерационная матрица: $\mathbf{B}(\tau) = \mathbf{I} - \tau\mathbf{A}$.

Метод сходится тогда и только тогда, когда спектральный радиус $\rho(\mathbf{B}(\tau)) < 1$.

Найдём собственные значения матрицы $\mathbf{A}$ из характеристического уравнения $\det(\mathbf{A} - \lambda\mathbf{I}) = 0$:
\begin{equation}
\det\begin{pmatrix} 3-\lambda & -1 & 1 \\ 2 & 6-\lambda & -2 \\ -1 & 1 & 9-\lambda \end{pmatrix} = 0.
\end{equation}

Раскрывая определитель:
\begin{align}
&(3-\lambda)[(6-\lambda)(9-\lambda) + 2] + [2(9-\lambda) - 2] + [2 + (6-\lambda)] \nonumber \\
&= (3-\lambda)(\lambda^2 - 15\lambda + 56) + (16 - 2\lambda) + (8 - \lambda) \nonumber \\
&= -\lambda^3 + 18\lambda^2 - 104\lambda + 192 = 0.
\end{align}

Проверяем $\lambda = 4$: $-64 + 288 - 416 + 192 = 0$ 

Делим на $(\lambda - 4)$: $\lambda^3 - 18\lambda^2 + 104\lambda - 192 = (\lambda - 4)(\lambda^2 - 14\lambda + 48) = 0$.

Корни квадратного уравнения: $\lambda = \frac{14 \pm \sqrt{196 - 192}}{2} = \frac{14 \pm 2}{2}$, т.е. $\lambda = 8$ и $\lambda = 6$.

\textbf{Собственные значения матрицы $\mathbf{A}$:} $\lambda_1 = 4$, $\lambda_2 = 6$, $\lambda_3 = 8$.

Собственные значения итерационной матрицы $\mathbf{B}(\tau) = \mathbf{I} - \tau\mathbf{A}$:
\begin{equation}
\mu_1 = 1 - 4\tau, \quad \mu_2 = 1 - 6\tau, \quad \mu_3 = 1 - 8\tau.
\end{equation}

\subsection*{а) Сходимость при бесконечной точности}

Метод сходится при $\rho(\mathbf{B}(\tau)) = \max(|\mu_1|, |\mu_2|, |\mu_3|) < 1$:
\begin{align}
|1 - 4\tau| &< 1 \implies 0 < \tau < \frac{1}{2}, \\
|1 - 6\tau| &< 1 \implies 0 < \tau < \frac{1}{3}, \\
|1 - 8\tau| &< 1 \implies 0 < \tau < \frac{1}{4}.
\end{align}

Пересечение условий:
\begin{equation}
\boxed{0 < \tau < \frac{1}{4}.}
\end{equation}

\subsection*{б) Сходимость при длине мантиссы 52 бита}

При конечной точности вычислений (double precision, $\varepsilon_{\text{маш}} = 2^{-52} \approx 2.22 \cdot 10^{-16}$) метод сходится к окрестности точного решения, если $\rho(\mathbf{B}(\tau)) < 1$. Однако при $\tau$, близком к границам интервала $(0, 1/4)$, спектральный радиус близок к 1, и сходимость становится чрезвычайно медленной.

Практически метод сходится за разумное число итераций, если $\rho(\mathbf{B}(\tau)) < 1 - c\varepsilon_{\text{маш}}$ для некоторой константы $c$. Это исключает малые окрестности точек $\tau = 0$ и $\tau = 1/4$.

Формально условие сходимости остаётся тем же:
\begin{equation}
\boxed{0 < \tau < \frac{1}{4},}
\end{equation}
но практически следует выбирать $\tau$ подальше от границ, например $\tau \in (10^{-15}, 1/4 - 10^{-15})$.

\subsection*{в) Оптимальный параметр для максимальной скорости сходимости}

Нужно минимизировать $\rho(\mathbf{B}(\tau)) = \max(|1-4\tau|, |1-6\tau|, |1-8\tau|)$.

При $\tau \in (0, 1/4)$:
\begin{itemize}
\item $1 - 4\tau$ убывает от 1 до 0,
\item $1 - 6\tau$ убывает от 1 до $-1/2$,
\item $1 - 8\tau$ убывает от 1 до $-1$.
\end{itemize}

Максимум по модулю определяется крайними собственными значениями: $\max(|1-4\tau|, |1-8\tau|)$.

Оптимум достигается, когда $|1-4\tau| = |1-8\tau|$:
\begin{equation}
1 - 4\tau = -(1 - 8\tau) \implies 2 = 12\tau \implies \tau_{\text{opt}} = \frac{1}{6}.
\end{equation}

При $\tau = 1/6$:
\begin{align}
\mu_1 &= 1 - 4 \cdot \frac{1}{6} = \frac{1}{3}, \\
\mu_2 &= 1 - 6 \cdot \frac{1}{6} = 0, \\
\mu_3 &= 1 - 8 \cdot \frac{1}{6} = -\frac{1}{3}.
\end{align}

\begin{equation}
\boxed{\tau_{\text{opt}} = \frac{1}{6}, \quad \rho(\mathbf{B}(\tau_{\text{opt}})) = \frac{1}{3}.}
\end{equation}

\subsection*{г) Оптимальный параметр для первого шага метода наискорейшего спуска}

При $\mathbf{x}^{(0)} = \mathbf{0}$: невязка $\mathbf{r}^{(0)} = \mathbf{f} - \mathbf{A}\mathbf{x}^{(0)} = \mathbf{f} = (4, -4, 2)^T$.

Метод наискорейшего спуска: $\mathbf{x}^{(1)} = \mathbf{x}^{(0)} + \tau_0 \mathbf{r}^{(0)} = \tau_0 \mathbf{f}$.

Оптимальный $\tau_0$ минимизирует $\|\mathbf{r}^{(1)}\|^2 = \|\mathbf{f} - \mathbf{A}\mathbf{x}^{(1)}\|^2 = \|\mathbf{f} - \tau_0 \mathbf{A}\mathbf{f}\|^2$.

Условие минимума: $\frac{d}{d\tau_0}\|\mathbf{f} - \tau_0 \mathbf{A}\mathbf{f}\|^2 = -2(\mathbf{f} - \tau_0 \mathbf{A}\mathbf{f}, \mathbf{A}\mathbf{f}) = 0$.

Отсюда:
\begin{equation}
\tau_0 = \frac{(\mathbf{f}, \mathbf{A}\mathbf{f})}{(\mathbf{A}\mathbf{f}, \mathbf{A}\mathbf{f})} = \frac{(\mathbf{f}, \mathbf{A}\mathbf{f})}{\|\mathbf{A}\mathbf{f}\|^2}.
\end{equation}

Вычислим $\mathbf{A}\mathbf{f}$:
\begin{equation}
\mathbf{A}\mathbf{f} = \begin{pmatrix} 3 & -1 & 1 \\ 2 & 6 & -2 \\ -1 & 1 & 9 \end{pmatrix} \begin{pmatrix} 4 \\ -4 \\ 2 \end{pmatrix} = \begin{pmatrix} 12 + 4 + 2 \\ 8 - 24 - 4 \\ -4 - 4 + 18 \end{pmatrix} = \begin{pmatrix} 18 \\ -20 \\ 10 \end{pmatrix}.
\end{equation}

Скалярное произведение $(\mathbf{f}, \mathbf{A}\mathbf{f})$:
\begin{equation}
(\mathbf{f}, \mathbf{A}\mathbf{f}) = 4 \cdot 18 + (-4) \cdot (-20) + 2 \cdot 10 = 72 + 80 + 20 = 172.
\end{equation}

Квадрат нормы $\|\mathbf{A}\mathbf{f}\|^2$:
\begin{equation}
\|\mathbf{A}\mathbf{f}\|^2 = 18^2 + (-20)^2 + 10^2 = 324 + 400 + 100 = 824.
\end{equation}

Оптимальный параметр:
\begin{equation}
\boxed{\tau_0 = \frac{172}{824} = \frac{43}{206} \approx 0.2087.}
\end{equation}

\subsection*{Итог}

\begin{itemize}
\item а) Метод сходится при $0 < \tau < 1/4$.
\item б) При конечной точности формально условие то же, но практически следует избегать $\tau$ близких к 0 и $1/4$.
\item в) Оптимальный параметр $\tau_{\text{opt}} = 1/6$, спектральный радиус $\rho = 1/3$.
\item г) Для первого шага метода наискорейшего спуска $\tau_0 = 43/206 \approx 0.2087$.
\end{itemize}








\section*{Задача III.4.3}

\textbf{Условие:} Предложить алгоритм проведения на плоскости окружности через четыре и более точек методом наименьших квадратов.

\vspace{0.5cm}
\textbf{Решение:}

\subsection*{1. Постановка задачи}

Пусть на плоскости заданы $N$ точек с координатами $(x_i, y_i)$, $i = 1, 2, \ldots, N$, причём $N \geqslant 4$. Требуется найти окружность, которая наилучшим образом (в смысле метода наименьших квадратов) аппроксимирует эти точки.

Уравнение окружности с центром в точке $(a, b)$ и радиусом $R$ имеет вид:
\begin{equation}
(x - a)^2 + (y - b)^2 = R^2.
\end{equation}

Раскроем скобки:
\begin{equation}
x^2 - 2ax + a^2 + y^2 - 2by + b^2 = R^2.
\end{equation}

Перегруппируем слагаемые, перенеся всё, что зависит от координат точки, в правую часть:
\begin{equation}
2ax + 2by + (R^2 - a^2 - b^2) = x^2 + y^2.
\end{equation}

Введём обозначение $c = R^2 - a^2 - b^2$. Тогда уравнение окружности принимает вид:
\begin{equation}
2ax + 2by + c = x^2 + y^2.
\label{eq:43_lin}
\end{equation}

\textbf{Ключевое наблюдение:} уравнение~\eqref{eq:43_lin} линейно относительно неизвестных параметров $a$, $b$ и $c$. Это позволяет свести задачу к решению переопределённой системы линейных алгебраических уравнений (СЛАУ).

\subsection*{2. Составление переопределённой СЛАУ}

Подставим координаты каждой из $N$ заданных точек в уравнение~\eqref{eq:43_lin}. Получим систему:
\begin{equation}
\begin{cases}
2a x_1 + 2b y_1 + c = x_1^2 + y_1^2, \\
2a x_2 + 2b y_2 + c = x_2^2 + y_2^2, \\
\quad \vdots \\
2a x_N + 2b y_N + c = x_N^2 + y_N^2.
\end{cases}
\end{equation}

В матричной форме эта система записывается как:
\begin{equation}
\mathbf{G} \boldsymbol{\theta} = \mathbf{f},
\label{eq:43_slau}
\end{equation}
где
\begin{equation}
\mathbf{G} =
\begin{pmatrix}
2x_1 & 2y_1 & 1 \\
2x_2 & 2y_2 & 1 \\
\vdots & \vdots & \vdots \\
2x_N & 2y_N & 1
\end{pmatrix}, \quad
\boldsymbol{\theta} =
\begin{pmatrix}
a \\ b \\ c
\end{pmatrix}, \quad
\mathbf{f} =
\begin{pmatrix}
x_1^2 + y_1^2 \\
x_2^2 + y_2^2 \\
\vdots \\
x_N^2 + y_N^2
\end{pmatrix}.
\end{equation}

Поскольку $N \geqslant 4 > 3$, система~\eqref{eq:43_slau} является переопределённой и, вообще говоря, не имеет точного решения.

\subsection*{3. Метод наименьших квадратов}

Вместо точного решения будем искать вектор $\boldsymbol{\theta}$, минимизирующий квадрат евклидовой нормы невязки:
\begin{equation}
\Phi(\boldsymbol{\theta}) = \|\mathbf{G}\boldsymbol{\theta} - \mathbf{f}\|_2^2 = \sum_{i=1}^{N} \bigl(2a x_i + 2b y_i + c - (x_i^2 + y_i^2)\bigr)^2 \to \min.
\end{equation}

Необходимое условие минимума --- обращение в нуль градиента $\Phi$ по всем компонентам $\boldsymbol{\theta}$:
\begin{equation}
\frac{\partial \Phi}{\partial \boldsymbol{\theta}} = 2\mathbf{G}^T(\mathbf{G}\boldsymbol{\theta} - \mathbf{f}) = \mathbf{0}.
\end{equation}

Отсюда получаем \textbf{нормальную систему} уравнений:
\begin{equation}
\mathbf{G}^T \mathbf{G} \boldsymbol{\theta} = \mathbf{G}^T \mathbf{f}.
\label{eq:43_normal}
\end{equation}

Матрица $\mathbf{G}^T \mathbf{G}$ имеет размер $3 \times 3$. Если точки не лежат на одной прямой, эта матрица невырождена, и решение единственно:
\begin{equation}
\boldsymbol{\theta} = (\mathbf{G}^T \mathbf{G})^{-1} \mathbf{G}^T \mathbf{f}.
\label{eq:43_sol}
\end{equation}

\subsection*{4. Восстановление радиуса}

После того как найдены параметры $a$, $b$, $c$, радиус окружности восстанавливается из соотношения $c = R^2 - a^2 - b^2$:
\begin{equation}
R = \sqrt{a^2 + b^2 + c}.
\end{equation}

\subsection*{5. Итоговый алгоритм}

\begin{enumerate}
\item \textbf{Входные данные:} координаты $N$ точек $(x_i, y_i)$, $i = 1, \ldots, N$, $N \geqslant 4$.

\item \textbf{Шаг 1.} Сформировать матрицу $\mathbf{G}$ размера $N \times 3$ и вектор $\mathbf{f}$ длины $N$:
\begin{equation}
G_{i1} = 2x_i, \quad G_{i2} = 2y_i, \quad G_{i3} = 1, \quad f_i = x_i^2 + y_i^2.
\end{equation}

\item \textbf{Шаг 2.} Вычислить матрицу $\mathbf{A} = \mathbf{G}^T \mathbf{G}$ размера $3 \times 3$ и вектор $\mathbf{b} = \mathbf{G}^T \mathbf{f}$ длины $3$.

\item \textbf{Шаг 3.} Решить нормальную систему $\mathbf{A} \boldsymbol{\theta} = \mathbf{b}$ относительно $\boldsymbol{\theta} = (a, b, c)^T$.

\item \textbf{Шаг 4.} Вычислить радиус: $R = \sqrt{a^2 + b^2 + c}$.

\item \textbf{Выходные данные:} центр окружности $(a, b)$ и радиус $R$.
\end{enumerate}

\subsection*{6. Замечание о численной устойчивости}

На практике для повышения численной устойчивости алгоритма рекомендуется перед построением матрицы $\mathbf{G}$ выполнить центрирование данных: вычесть из всех координат $x_i$ и $y_i$ их средние значения $\bar{x}$ и $\bar{y}$. После нахождения параметров в центрированной системе координат центр окружности пересчитывается обратно: $a_{\text{исх}} = a + \bar{x}$, $b_{\text{исх}} = b + \bar{y}$.

\newpage



\section*{Задача III.4.4}

\textbf{Условие:} Напряжённость магнитного поля $H$ и магнитная индукция $B$ связаны соотношением $B = H/(a + bH)$. По результатам экспериментальных измерений определить коэффициенты $a$ и $b$:

\begin{center}
\begin{tabular}{|c|c|c|c|c|c|c|c|c|}
\hline
$H$ & 8 & 10 & 15 & 20 & 30 & 40 & 60 & 80 \\
\hline
$B$ & 13,0 & 14,0 & 15,4 & 16,3 & 17,2 & 17,8 & 18,5 & 18,8 \\
\hline
\end{tabular}
\end{center}

\vspace{0.5cm}
\textbf{Решение:}

\subsection*{1. Теоретическое обоснование метода итерированного веса}

Исходная задача состоит в минимизации нелинейного функционала невязки:
\begin{equation}
\Phi(a, b) = \sum_{i=1}^{N} \left(B_i - \frac{H_i}{a + bH_i}\right)^2 \to \min.
\label{eq:nonlinear_func}
\end{equation}

Прямая минимизация~\eqref{eq:nonlinear_func} приводит к системе нелинейных уравнений. Чтобы свести задачу к линейной, выполним алгебраическое преобразование \textit{внутри} квадрата невязки. Приведём выражение в скобках к общему знаменателю:
\begin{equation}
B_i - \frac{H_i}{a + bH_i} = \frac{aB_i + bB_iH_i - H_i}{a + bH_i}.
\end{equation}

Подставим это обратно в функционал~\eqref{eq:nonlinear_func}:
\begin{equation}
\Phi(a, b) = \sum_{i=1}^{N} \frac{(aB_i + bB_iH_i - H_i)^2}{(a + bH_i)^2} = \sum_{i=1}^{N} w_i(a, b) \cdot (aB_i + bB_iH_i - H_i)^2 \to \min,
\label{eq:weighted_func}
\end{equation}
где введены \textbf{веса}:
\begin{equation}
w_i(a, b) = \frac{1}{(a + bH_i)^2}.
\end{equation}

Формула~\eqref{eq:weighted_func} представляет собой задачу \textbf{взвешенного метода наименьших квадратов (ВМНК)} для линейного уравнения $aB_i + bB_iH_i = H_i$. Однако веса $w_i$ зависят от искомых параметров $a$ и $b$, поэтому мы не можем решить эту систему за один шаг. Для этого применяется \textbf{метод итерированного веса} (Iteratively Reweighted Least Squares, IRLS).

\subsection*{2. Алгоритм метода итерированного веса}

\begin{enumerate}
    \item \textbf{Нулевая итерация (начальное приближение):} Решаем задачу \textit{невзвешенного} МНК (полагая $w_i^{(0)} = 1$), чтобы получить начальные оценки $a_0$ и $b_0$.
    \item \textbf{$k$-я итерация:} 
    \begin{enumerate}
        \item Вычисляем веса на основе предыдущих оценок: $w_i^{(k)} = \frac{1}{(a_{k-1} + b_{k-1}H_i)^2}$.
        \item Решаем взвешенную нормальную систему:
        \begin{equation}
        \begin{cases}
        a \sum w_i^{(k)} B_i^2 + b \sum w_i^{(k)} B_i^2 H_i = \sum w_i^{(k)} B_i H_i, \\
        a \sum w_i^{(k)} B_i^2 H_i + b \sum w_i^{(k)} B_i^2 H_i^2 = \sum w_i^{(k)} B_i H_i^2.
        \end{cases}
        \end{equation}
        \item Получаем новые оценки $a_k$ и $b_k$.
    \end{enumerate}
    \item \textbf{Критерий остановки:} Повторяем шаг 2 до тех пор, пока изменения параметров $|a_k - a_{k-1}|$ и $|b_k - b_{k-1}|$ не станут меньше заданной точности (например, $10^{-4}$).
\end{enumerate}

\subsection*{3. Численная реализация}

\textbf{Шаг 0. Невзвешенное начальное приближение.}
Решая систему с $w_i = 1$ (как в классическом линейном МНК), получаем:
\begin{equation}
a_0 \approx 0{,}218, \quad b_0 \approx 0{,}0506.
\end{equation}

\textbf{Шаг 1. Первая итерация с весами.}
Вычислим веса $w_i^{(1)} = \frac{1}{(0{,}218 + 0{,}0506 \cdot H_i)^2}$ для каждой точки:
\begin{center}
\begin{tabular}{|c|c|c|c|}
\hline
$H_i$ & $a_0 + b_0 H_i$ & $w_i^{(1)}$ & $w_i^{(1)} B_i^2$ \\
\hline
8 & 0,623 & 2,57 & 434,3 \\
10 & 0,724 & 1,91 & 374,4 \\
15 & 0,977 & 1,05 & 248,9 \\
20 & 1,230 & 0,66 & 174,9 \\
30 & 1,736 & 0,33 & 97,4 \\
40 & 2,242 & 0,20 & 63,2 \\
60 & 3,254 & 0,09 & 30,8 \\
80 & 4,266 & 0,05 & 17,7 \\
\hline
\end{tabular}
\end{center}

Подставляя эти веса в нормальную систему, вычисляем суммы:
\begin{align*}
\sum w_i B_i^2 &\approx 1441,6, & \sum w_i B_i^2 H_i &\approx 23177, \\
\sum w_i B_i^2 H_i^2 &\approx 606100, & \sum w_i B_i H_i &\approx 1477, \\
\sum w_i B_i H_i^2 &\approx 35456. & &
\end{align*}

Решаем систему:
\begin{equation}
\begin{pmatrix}
1441{,}6 & 23177 \\
23177 & 606100
\end{pmatrix}
\begin{pmatrix}
a_1 \\ b_1
\end{pmatrix}
=
\begin{pmatrix}
1477 \\ 35456
\end{pmatrix}.
\end{equation}

Определитель системы $\Delta \approx 3{,}36 \cdot 10^8$. Находим:
\begin{equation}
a_1 = \frac{1477 \cdot 606100 - 23177 \cdot 35456}{\Delta} \approx 0{,}218, \quad
b_1 = \frac{1441{,}6 \cdot 35456 - 1477 \cdot 23177}{\Delta} \approx 0{,}0501.
\end{equation}

\textbf{Шаг 2 и далее.} 
Повторение процедуры с новыми весами $w_i^{(2)} = \frac{1}{(0{,}218 + 0{,}0501 \cdot H_i)^2}$ показывает, что значения $a$ и $b$ практически перестают меняться (изменения меньше $10^{-4}$). Метод сошёлся.

\subsection*{4. Итоговые результаты и проверка}

Оптимальные значения коэффициентов:
\begin{equation}
\boxed{a \approx 0{,}218, \quad b \approx 0{,}0501.}
\end{equation}

Уравнение кривой намагничивания:
\begin{equation}
B = \frac{H}{0{,}218 + 0{,}0501H}.
\end{equation}

Проверим качество аппроксимации, вычислив невязки $\Delta B_i = B_{i,\text{изм}} - B_{i,\text{расч}}$:

\begin{center}
\begin{tabular}{|c|c|c|c|}
\hline
$H$ & $B_{\text{изм}}$ & $B_{\text{расч}}$ & Невязка $\Delta B$ \\
\hline
8 & 13,0 & 12,93 & +0,07 \\
10 & 14,0 & 13,88 & +0,12 \\
15 & 15,4 & 15,38 & +0,02 \\
20 & 16,3 & 16,30 & 0,00 \\
30 & 17,2 & 17,31 & -0,11 \\
40 & 17,8 & 17,86 & -0,06 \\
60 & 18,5 & 18,46 & +0,04 \\
80 & 18,8 & 18,77 & +0,03 \\
\hline
\end{tabular}
\end{center}

Сумма невязок близка к нулю ($\sum \Delta B_i \approx 0{,}11$), а максимальное отклонение не превышает $0{,}12$ (менее $1\%$ от максимального значения $B$).

\subsection*{5. Физическая интерпретация}
\begin{itemize}
    \item Параметр $a \approx 0{,}218$ определяет начальную магнитную проницаемость: $\mu_0 \approx 1/a \approx 4{,}6$.
    \item Параметр $b \approx 0{,}0501$ определяет уровень магнитного насыщения: $B_{\text{нас}} = 1/b \approx 19{,}96$ Тл, что полностью согласуется с поведением ферромагнетиков при больших $H$.
\end{itemize}


\section*{Задача III.4.8}

\textbf{Условие:} Пусть замеры функции $y = f(x)$ осуществлены в точках $x_k = \cos\frac{\pi(1+2k)}{2(n+1)}$, $k = 0, 1, \ldots, n$, являющихся нулями многочлена Чебышёва $T_{n+1}(x)$, и собраны в таблицу. Среди многочленов степени не выше заданного $k$, $0 \le k \le n$, указать тот многочлен $P_k(x)$, который наилучшим (в смысле метода наименьших квадратов) образом приближает заданную функцию.

Провести вычисления в случае $n = 3$:

\begin{center}
\begin{tabular}{|c|c|c|c|c|}
\hline
$x$ & $x_0$ & $x_1$ & $x_2$ & $x_3$ \\
\hline
$y$ & 2 & 1 & 3 & 0 \\
\hline
\end{tabular}
\end{center}

\vspace{0.5cm}
\textbf{Решение:}

\subsection*{1. Теоретическая часть}

Ищем приближающий многочлен в виде разложения по многочленам Чебышёва:
\begin{equation}
P_k(x) = \sum_{j=0}^{k} C_j T_j(x), \quad 0 \le k \le n.
\label{eq:48_pk}
\end{equation}

Функционал метода наименьших квадратов имеет вид:
\begin{equation}
\Phi(C_0, C_1, \ldots, C_k) = \sum_{m=0}^{n} \left(y_m - \sum_{j=0}^{k} C_j T_j(x_m)\right)^2 \to \min.
\end{equation}

Ключевое свойство, на которое указывает условие задачи: многочлены Чебышёва $T_j(x)$ образуют \textbf{ортогональную систему} на множестве нулей $T_{n+1}(x)$ с единичными весами. А именно, для $x_m = \cos\frac{\pi(1+2m)}{2(n+1)}$ выполняется:
\begin{equation}
\sum_{m=0}^{n} T_i(x_m) T_j(x_m) = 
\begin{cases} 
0, & i \neq j, \\
n+1, & i = j = 0, \\
\frac{n+1}{2}, & i = j \neq 0.
\end{cases}
\label{eq:48_orth}
\end{equation}

Благодаря ортогональности~\eqref{eq:48_orth} нормальная система МНК \textbf{диагонализируется}, и каждый коэффициент $C_j$ находится независимо из условия $\frac{\partial \Phi}{\partial C_j} = 0$:
\begin{equation}
C_j = \frac{\displaystyle\sum_{m=0}^{n} y_m T_j(x_m)}{\displaystyle\sum_{m=0}^{n} T_j^2(x_m)}.
\label{eq:48_cj}
\end{equation}

Подставляя значения норм из~\eqref{eq:48_orth}, получаем явные формулы:
\begin{equation}
C_0 = \frac{1}{n+1} \sum_{m=0}^{n} y_m, \qquad C_j = \frac{2}{n+1} \sum_{m=0}^{n} y_m T_j(x_m), \quad j = 1, 2, \ldots, k.
\label{eq:48_cj_expl}
\end{equation}

\subsection*{2. Вычисления для $n = 3$}

Узлы интерполяции --- нули многочлена $T_4(x)$:
\begin{equation}
x_m = \cos\frac{\pi(1+2m)}{8}, \quad m = 0, 1, 2, 3.
\end{equation}

Выпишем их явно:
\begin{align}
x_0 &= \cos\frac{\pi}{8} = \frac{\sqrt{2+\sqrt{2}}}{2} \approx 0{,}92388, \\
x_1 &= \cos\frac{3\pi}{8} = \frac{\sqrt{2-\sqrt{2}}}{2} \approx 0{,}38268, \\
x_2 &= \cos\frac{5\pi}{8} = -\frac{\sqrt{2-\sqrt{2}}}{2} \approx -0{,}38268, \\
x_3 &= \cos\frac{7\pi}{8} = -\frac{\sqrt{2+\sqrt{2}}}{2} \approx -0{,}92388.
\end{align}

Значения многочленов Чебышёва $T_j(x_m) = \cos\frac{j\pi(1+2m)}{8}$ в узлах:

\begin{center}
\begin{tabular}{|c|c|c|c|c|}
\hline
 & $m=0$ & $m=1$ & $m=2$ & $m=3$ \\
\hline
$T_0(x_m)$ & 1 & 1 & 1 & 1 \\
\hline
$T_1(x_m)$ & $\cos\frac{\pi}{8}$ & $\cos\frac{3\pi}{8}$ & $-\cos\frac{3\pi}{8}$ & $-\cos\frac{\pi}{8}$ \\
\hline
$T_2(x_m)$ & $\cos\frac{\pi}{4} = \frac{\sqrt{2}}{2}$ & $\cos\frac{3\pi}{4} = -\frac{\sqrt{2}}{2}$ & $\cos\frac{5\pi}{4} = -\frac{\sqrt{2}}{2}$ & $\cos\frac{7\pi}{4} = \frac{\sqrt{2}}{2}$ \\
\hline
$T_3(x_m)$ & $\cos\frac{3\pi}{8}$ & $-\cos\frac{\pi}{8}$ & $\cos\frac{\pi}{8}$ & $-\cos\frac{3\pi}{8}$ \\
\hline
\end{tabular}
\end{center}

\textbf{Коэффициент $C_0$:}
\begin{equation}
C_0 = \frac{1}{4}(y_0 + y_1 + y_2 + y_3) = \frac{1}{4}(2 + 1 + 3 + 0) = \frac{3}{2}.
\end{equation}

\textbf{Коэффициент $C_1$:}
\begin{align}
C_1 &= \frac{2}{4} \sum_{m=0}^{3} y_m T_1(x_m) = \frac{1}{2}\left(2\cos\frac{\pi}{8} + 1\cdot\cos\frac{3\pi}{8} + 3\cdot\left(-\cos\frac{3\pi}{8}\right) + 0\cdot\left(-\cos\frac{\pi}{8}\right)\right) \nonumber \\
&= \frac{1}{2}\left(2\cos\frac{\pi}{8} - 2\cos\frac{3\pi}{8}\right) = \cos\frac{\pi}{8} - \cos\frac{3\pi}{8} = \frac{\sqrt{2+\sqrt{2}} - \sqrt{2-\sqrt{2}}}{2} \approx 0{,}5412.
\end{align}

\textbf{Коэффициент $C_2$:}
\begin{align}
C_2 &= \frac{2}{4} \sum_{m=0}^{3} y_m T_2(x_m) = \frac{1}{2}\left(2\cdot\frac{\sqrt{2}}{2} + 1\cdot\left(-\frac{\sqrt{2}}{2}\right) + 3\cdot\left(-\frac{\sqrt{2}}{2}\right) + 0\cdot\frac{\sqrt{2}}{2}\right) \nonumber \\
&= \frac{1}{2}\left(\sqrt{2} - \frac{\sqrt{2}}{2} - \frac{3\sqrt{2}}{2}\right) = \frac{1}{2}\left(-\sqrt{2}\right) = -\frac{\sqrt{2}}{2} \approx -0{,}7071.
\end{align}

\textbf{Коэффициент $C_3$:}
\begin{align}
C_3 &= \frac{2}{4} \sum_{m=0}^{3} y_m T_3(x_m) = \frac{1}{2}\left(2\cos\frac{3\pi}{8} + 1\cdot\left(-\cos\frac{\pi}{8}\right) + 3\cdot\cos\frac{\pi}{8} + 0\cdot\left(-\cos\frac{3\pi}{8}\right)\right) \nonumber \\
&= \frac{1}{2}\left(2\cos\frac{3\pi}{8} + 2\cos\frac{\pi}{8}\right) = \cos\frac{\pi}{8} + \cos\frac{3\pi}{8} = \frac{\sqrt{2+\sqrt{2}} + \sqrt{2-\sqrt{2}}}{2} \approx 1{,}3066.
\end{align}

\subsection*{3. Искомые многочлены наилучшего приближения}

Подставляя найденные коэффициенты в формулу~\eqref{eq:48_pk}, получаем:

\begin{align}
P_0(x) &= \frac{3}{2}, \\
P_1(x) &= \frac{3}{2} + \left(\cos\frac{\pi}{8} - \cos\frac{3\pi}{8}\right) T_1(x) = \frac{3}{2} + \left(\cos\frac{\pi}{8} - \cos\frac{3\pi}{8}\right) x, \\
P_2(x) &= P_1(x) - \frac{\sqrt{2}}{2}\, T_2(x) = \frac{3}{2} + \left(\cos\frac{\pi}{8} - \cos\frac{3\pi}{8}\right) x - \frac{\sqrt{2}}{2}(2x^2 - 1), \\
P_3(x) &= P_2(x) + \left(\cos\frac{\pi}{8} + \cos\frac{3\pi}{8}\right) T_3(x) \nonumber \\
&= \frac{3}{2} + \left(\cos\frac{\pi}{8} - \cos\frac{3\pi}{8}\right) x - \frac{\sqrt{2}}{2}(2x^2 - 1) + \left(\cos\frac{\pi}{8} + \cos\frac{3\pi}{8}\right)(4x^3 - 3x).
\end{align}

\textbf{Замечание.} Поскольку при $n = 3$ многочлен $P_3(x)$ имеет степень 3 и строится по 4 узлам, он совпадает с интерполяционным многочленом Лагранжа и точно проходит через все заданные точки. Многочлены $P_0$, $P_1$, $P_2$ дают наилучшее среднеквадратичное приближение среди всех многочленов соответствующей степени.

\subsection*{4. Численная проверка}

Для контроля вычислим значения $P_3(x_m)$ и убедимся, что они совпадают с табличными $y_m$. Например, для $m = 0$:
\begin{align}
P_3(x_0) &= C_0 + C_1 T_1(x_0) + C_2 T_2(x_0) + C_3 T_3(x_0) \nonumber \\
&= \frac{3}{2} + \left(\cos\frac{\pi}{8} - \cos\frac{3\pi}{8}\right)\cos\frac{\pi}{8} - \frac{\sqrt{2}}{2}\cdot\frac{\sqrt{2}}{2} + \left(\cos\frac{\pi}{8} + \cos\frac{3\pi}{8}\right)\cos\frac{3\pi}{8} \nonumber \\
&= \frac{3}{2} + \cos^2\frac{\pi}{8} - \cos\frac{\pi}{8}\cos\frac{3\pi}{8} - \frac{1}{2} + \cos\frac{\pi}{8}\cos\frac{3\pi}{8} + \cos^2\frac{3\pi}{8} \nonumber \\
&= 1 + \cos^2\frac{\pi}{8} + \cos^2\frac{3\pi}{8} = 1 + \frac{1+\cos\frac{\pi}{4}}{2} + \frac{1+\cos\frac{3\pi}{4}}{2} = 1 + \frac{1}{2} + \frac{1}{2} = 2 = y_0. \quad \checkmark
\end{align}

Аналогично проверяются остальные узлы.




\end{document}